\section{Related Work}
\label{sec:relatedwork}
\paragraph{Python type checking and type-error detection.}
Python's annotation system supplies optional type declarations, and
structural protocols describe behavior shared by different
classes~\cite{pep484,pep544}. Checkers combine this interface information
with inference and local narrowing~\cite{mypyTypeNarrowing,pythonTypeNarrowing}.
Empirical studies examine how gradual typing affects checker behavior
and which types are needed to describe dynamic
objects~\cite{xu2023gradual,sun2023dynamicObjects}. Our analysis focuses
on the state at an operation after preceding calls have changed shared
locations. The relevant relation connects a calling condition to the
particular update that reaches that operation.

Pyinder is a close task-level comparison. It infers types, analyzes
functions under parameter-type contexts, and aggregates and filters member
types to identify Python type errors~\cite{oh2024pyinder}. We study a
representation that keeps entry state, execution conditions, ordered
requirements, and caller-visible updates together. Actual object bindings
instantiate this representation at calls, while conflict dependencies
identify the callee type paths used for re-inference. The whole-tool
comparison measures detection outcomes; the controlled variants examine
the roles of relation preservation and feedback.

\paragraph{Type inference and generated annotations.}
Annotation inference reconstructs interface types from code and repository
context. PyTIR models entity dependencies and iteratively updates types
and dependencies with checker feedback~\cite{sun2025pytir}. Its evaluation
includes annotation accuracy and diagnostics introduced by generated
annotations. Our target is the original program's operation behavior and
existing annotation contracts. The state analysis uses a fixed pyflow
call graph~\cite{pyflow2026}; refinement changes type-state relations and
their dependent judgments. The inference-plus-checker baseline evaluates
how annotation generation serves this detection task, separating original
errors from generated-interface inconsistencies.

\paragraph{Relational and compositional program analysis.}
Hoare logic connects assertions at program entry and exit~\cite{hoare1969axiomatic}. Compositional symbolic execution
reuses procedure behavior at calls~\cite{godefroid2007compositional}, while
abstract interpretation provides a foundation for reasoning about program
states and their relations~\cite{cousot1977abstract}. MOPSA applies abstract
interpretation to Python type analysis, including state relations and
control and exception behavior~\cite{monat2020python}.

Our model uses these relational and compositional principles for Python
type-error detection. Its cases preserve ordered checks and writes along
with entry and exit assertions, so intermediate failure dependencies remain
available during composition and refinement. The accompanying construction
procedure shares entry-independent facts and proposes implicit type
alternatives; detection revisits disputed alternatives through selectively
expanded type paths. The comparison with failure-preserving relational
summaries and direct body analysis evaluates these representation and
construction choices under matched rules.
